\begin{resumo}[RESUMO]
\vspace*{-6mm}

Sistemas de controle industrial de grande porte operam com centenas ou milhares de malhas de controle realimentadas, cuja sintonização inadequada pode causar oscilações, desperdício de energia e acionamento de intertravamentos de segurança. O diagnóstico automático da qualidade dessas malhas requer um ambiente de experimentação seguro e reprodutível, no qual distúrbios possam ser introduzidos e estratégias supervisórias testadas sem risco à planta real. Este trabalho propõe e implementa um laboratório de \textbf{gêmeo digital} do Processo Tennessee Eastman (TEP) — \textit{benchmark} clássico da comunidade de controle de processos — com foco em \textbf{integração de sistemas de controle}, demonstrando como simuladores, interfaces, orquestradores e controladores podem interoperar através de padrões abertos e arquiteturas modernas. A planta química é simulada em Rust com integração numérica Runge-Kutta de 4ª ordem e exposta via servidor gRPC; uma Interface Homem-Máquina (IHM) em FastAPI com WebSocket permite monitoramento em tempo real; e um Operator Kubernetes implementado em Go com Kubebuilder observa continuamente o estado da planta por meio de um Recurso Customizado (\textit{Custom Resource Definition} — CRD) denominado \texttt{PLCMachine}. Dezessete experimentos foram conduzidos, caracterizando o comportamento da planta no estado estacionário e sob os distúrbios IDV(1), IDV(2) e IDV(3). Os resultados demonstram que IDV(1) colapsa a planta por sobrepressão em aproximadamente 2,5 horas simuladas e IDV(2) causa colapso lento por acúmulo de inerte em dezenas de horas — ambos casos relevantes para o desenvolvimento de lógica supervisória. O Operator Kubernetes demonstrou conectividade estável com a planta e operação contínua em modo de observação, validando a arquitetura como infraestrutura para futura incorporação de métricas de qualidade de malha e integração via OPC UA.

\textbf{Palavras-chave}: Gêmeo digital. Integração de sistemas de controle. Processo Tennessee Eastman. Kubernetes. Operator pattern. Supervisão industrial. gRPC. Interoperabilidade.
\end{resumo}
